\documentclass[10pt,a4paper]{amsart}
\usepackage[margin=19mm]{geometry}
\usepackage{amsmath,amssymb,mathtools}
\usepackage[scr=rsfs,cal=euler]{mathalfa}
\usepackage{xfrac}
\usepackage[unicode,hidelinks,bookmarks=false]{hyperref}
\newcommand{\ie}{i.e.\ }
\newcommand{\cf}{cf.\ }
\newcommand{\resp}{respectively\ }
\newcommand{\Subsection}{\subsection}
\providecommand{\nobreakdash}{\nobreak\hbox{-}}
\setlength{\emergencystretch}{2em}
\multlinegap0pt
\usepackage{bm}
\usepackage{bbm}
\usepackage{dsfont}
\usepackage{xspace}
\usepackage{cancel}
\usepackage{mathtools}
\usepackage[shortlabels]{enumitem}
\usepackage{amsthm}
\makeatletter
\@ifundefined{define}{\newtheorem{define}{Define}}{}
\@ifundefined{lemma}{\newtheorem{lemma}[define]{Lemma}}{}
\@ifundefined{thm}{\newtheorem{thm}{Thm}}{}
\makeatother
\DeclareMathOperator{\diam}{diam}
\DeclareMathOperator{\proj}{proj}
\newcommand{\N}{\mathbb{N}}
\newcommand{\R}{\mathbb{R}}
\newcommand{\Z}{\mathbb{Z}}
\newcommand{\abs}[1]{\left|#1\right|}
\newcommand{\bilip}{\operatorname{bilip}}
\newcommand{\br}[1]{\left(#1\right)}
\newcommand{\cl}[1]{\overline{#1}}
\newcommand{\enorm}[1]{\left\|#1\right\|}
\newcommand{\mf}[1]{\mathfrak{#1}}
\newcommand{\set}[1]{\left\{#1\right\}}
\newcommand{\sqbr}[1]{\left[#1\right]}
\title[Extending Bilipschitz Mappings between Separated Nets]{Extending Bilipschitz Mappings between Separated Nets}
\author{Michael Dymond, Vojt\v{e}ch Kalu\v{z}a}
\date{Source: arXiv:2507.22007v3 (CC BY 4.0)}
\begin{document}
\maketitle

\noindent\emph{Source and licence.} Extending Bilipschitz Mappings between Separated Nets. Version arXiv:2507.22007v3 (CC BY 4.0), distributed under the Creative Commons Attribution 4.0 International licence, as recorded in the arXiv licence metadata for this exact version and shown on the abstract page https://arxiv.org/abs/2507.22007v3. Full rights notice: licenses/arxiv-2507.22007-CC-BY-4.0.txt.\par\smallskip

\noindent\textit{Editorial scope.} Counted unit: the Lemma whose source label is \texttt{lemma:inj\_rounding} in the article (source0.tex lines 681--693, source label \texttt{lemma:inj\_rounding}), delivered together with the article's own complete original proof of it (source0.tex lines 694--770, one \texttt{proof} environment of 77 lines). No mathematics is written, repaired or re-derived: statement and proof are the article's own text line for line. Required context retained uncounted: thm:simultaneous\_switching (lines 482--493). Every definition, display and label of those lines is retained unchanged. Omitted from this package and not redistributed: 1--481, 494--680, 771--997. Nothing inside a retained excerpt is omitted or altered: every retained body is delivered exactly as the article's lines are written, so no omission receipt of any kind accompanies this package. Citations: the retained excerpts carry no citation command at all, so no bibliography entry of the article is transcribed and no reference is redistributed. Reference adaptation: none was needed and none was made; every reference inside the delivered unit resolves inside it, so no unresolved reference or citation is produced. Printed slips kept literal: no printed word, symbol, hypothesis or number of the counted statement or of its proof was altered. Layout and class: the article's own class is \texttt{scrartcl} with packages including amsmath, amssymb, enumerate, amsthm, todonotes, mathtools, babel, fontenc, inputenc, lmodern, tikz, thmtools, thm-restate, stackrel; the wrapper is the lane's pinned \texttt{amsart} editorial wrapper at 10pt on A4, which restates the theorem-like environments the retained text needs and adds the article's own macro definitions that the retained text needs (\texttt{\textbackslash N}, \texttt{\textbackslash R}, \texttt{\textbackslash Z}, \texttt{\textbackslash abs}, \texttt{\textbackslash bilip}, \texttt{\textbackslash br}, \texttt{\textbackslash cl}, \texttt{\textbackslash diam}, \texttt{\textbackslash enorm}, \texttt{\textbackslash mf}, \texttt{\textbackslash proj}, \texttt{\textbackslash set}, \texttt{\textbackslash sqbr}), with extra wrapper packages (bm, bbm, dsfont, xspace, cancel, mathtools, enumitem). The delivered numbering is the unit's own. Assets: no retained span contains a figure environment, a graphics-inclusion command, a PDF-page inclusion, a table or a TikZ picture; no image file is copied into this package, the package declares no asset dependency and no image file is redistributed with it. Licence: version arXiv:2507.22007v3 of this article is distributed under the Creative Commons Attribution 4.0 International licence, as recorded in the arXiv licence metadata for this exact version and shown on the abstract page https://arxiv.org/abs/2507.22007v3. A full rights notice is delivered at licenses/arxiv-2507.22007-CC-BY-4.0.txt. Characters: every retained excerpt is pure ASCII, so the delivered engine, XeLaTeX with its default Latin Modern text font, reports no missing character.
\begin{thm}\label{thm:simultaneous_switching}
	Let $X\subseteq \R^{d}$. For each $x\in X$ let $y_{x}\in\R^{d}$, $r_{x}$ be a positive real number and $\mf{U}_{x}:=B([x,y_{x}],r_{x})$.  Suppose that $(\mf{U}_{x})_{x\in X}$ is pairwise disjoint, 
	\begin{equation*}
	\sup_{x\in X,\, y_{x}\neq x}\frac{r_{x}}{\enorm{y_{x}-x}}\leq \frac{1}{2}\qquad\text{and}\qquad\sup_{x\in X,\, y_{x}\neq x}\frac{\enorm{y_{x}-x}}{r_{x}}<\infty.
	\end{equation*}
	Then there is a bilipschitz mapping $\Gamma\colon \R^{d}\to \R^{d}$ such that
	\begin{enumerate}[(i)]
		\item $\Gamma(x)=y_{x}$ and $\Gamma(y_{x})=x$ for all $x\in X$.
		\item $\Gamma(p)=p$ for all $\displaystyle p\in \R^{d}\setminus \bigcup_{x\in X, y_{x}\neq x}\mf{U}_{x}$.
		\item $\displaystyle \bilip(\Gamma)\leq \max\set{1, \sup_{x\in X,\, y_{x}\neq x}\frac{4\enorm{y_{x}-x}^{2}}{r_{x}^{2}}}$.
	\end{enumerate}
\end{thm}

\begin{lemma}\label{lemma:inj_rounding}
	Let $d,H\in\N$, $d\geq2$, $0<s\leq\frac{1}{4}$, $N\geq(2\sqrt{d}/s)^3H^{\frac{1}{d-1}}$,
	\begin{equation*}
	X_{1},X_{2},\ldots,X_{H}\subseteq \R^{d-1}\times \sqbr{\frac{1}{2}+s,H+\frac{1}{2}-s} 
	\end{equation*}
	be pairwise disjoint and suppose that $X:=\displaystyle\bigcup_{m=1}^{H}X_{m}$ is $s$-separated. 
	Then there exists a $\displaystyle 2^{8}N^{2}H^{2}$-bilipschitz mapping $\Psi\colon \R^{d}\to\R^{d}$ such that
	\begin{enumerate}[(i)]
		\item\label{Xi_i} $\Psi(x)=x$ for all $x\in\R^{d-1}\times\br{\R\setminus \br{\frac{1}{2},H+\frac{1}{2}}}$.
		\item\label{Xi_ii} $\Psi(X_{m})\subseteq \br{\frac{1}{N}\Z^{d-1}\setminus \Z^{d-1}}\times\set{m}$ for $m=1,2,\ldots,H$.
		\item\label{Xi_iii} $\enorm{\proj_{\R^{d-1}}\circ \Psi(x)-\proj_{\R^{d-1}}(x)}\leq s^{2}$ for all $x\in X$.
	\end{enumerate}
\end{lemma}

\begin{proof}
	We will construct $\Psi$ as a composition $\Omega\circ \Xi$, where $\Xi$ will take care of the first $(d-1)$-coordinates and $\Omega$ of the last coordinate. 
	\paragraph{Construction of $\Xi$.}
	We will define a $16$-bilipschitz mapping $\Xi\colon\R^{d}\to\R^{d}$ with the following properties
	\begin{enumerate}[(A)]
		\item\label{Psi_i} $\Xi(x)=x$ for all $x\in\R^{d-1}\times \br{\R\setminus \br{\frac{1}{2},H+\frac{1}{2}}}$,
		\item\label{Psi_ii} $\enorm{\Xi(x)-x}\leq s^{2}$ for all $x\in X$. 
		\item\label{Psi_iii} $\proj_{\R^{d-1}}\circ \Xi(X)\subseteq \frac{1}{N}\Z^{d-1}\setminus \Z^{d-1}$ and $\proj_{\R^{d-1}}\circ \Xi|_{X}$ is injective.
		\item\label{Psi_iv} $\proj_{d}\circ \Xi(x)=\proj_{d}(x)$ for all $x\in X$.
	\end{enumerate}
	Let $Q:=\left[0,\frac{s^{2}}{\sqrt{d-1}}\right)^{d-1}$ and for each $x\in X$ let $\Phi(x)\in\Z^{d-1}$ be defined by the condition
	\begin{equation*}
	\proj_{\R^{d-1}}(x)\in Q+\frac{s^{2}}{\sqrt{d-1}}\Phi(x).
	\end{equation*}
	For each $z\in\Z^{d-1}$ we have that $\bigcup_{x\in\Phi^{-1}(\set{z})}B(x,s/2)$ is a pairwise disjoint union of balls all contained in the set $\br{\cl{B}_{\R^{d-1}}(Q,s/2)+\frac{s^{2}z}{\sqrt{d-1}}}\times\sqbr{\frac{1}{2},H+\frac{1}{2}}$. Hence, by comparing volumes and writing $\alpha_{d}$ for the $d$-dimensional volume of the unit ball in $\R^{d}$, we get
	\begin{equation}\label{eq:Phi_preim_upper_bound}
	\abs{\Phi^{-1}(\set{z})}\leq \frac{(s+s^{2})^{d-1}H}{\alpha_{d}\br{s/2}^{d}}\leq\frac{(2s)^{d-1}2^{d}H}{\alpha_{d}s^{d}}\leq \frac{2^{d-1}\sqrt{d^{d}}H}{s}.
	\end{equation}
	On the other hand we have that 
	\begin{equation}\label{eq:tricky}
	\abs{\br{Q+\frac{s^{2}}{\sqrt{d-1}}z}\cap \br{\frac{1}{N}\Z^{d-1}\setminus\Z^{d-1}}}\geq \br{\frac{s^{2}/\sqrt{d}}{1/N}-1}^{d-1}-1\geq \frac{1}{2}\br{\frac{Ns^{2}}{2\sqrt{d}}}^{d-1} .
	\end{equation}
	The lower bound on $N$ in the hypothesis is used for the last inequality. It also ensures that the final quantity of \eqref{eq:tricky} is at least that of \eqref{eq:Phi_preim_upper_bound}. Therefore, there exists for each $z\in\Z^{d-1}$ an injective mapping  
	\begin{equation*}
	\iota_{z}\colon \Phi^{-1}(\set{z})\to \br{Q+\frac{s^{2}}{\sqrt{d-1}}z}\cap \br{\frac{1}{N}\Z^{d-1}\setminus \Z^{d-1}}.
	\end{equation*}
	For each $x\in X$ and a radius $r_{x}=r_{x}(s)$ to be determined later, we let 
	\begin{equation*}
	y_{x}:=(\iota_{\Phi(x)}(x),x_{d})\in \R^{d}, \qquad \mf{U}_{x}:=B([x,y_{x}],r_{x}).
	\end{equation*}
	Observe that for each $x\in X$ the line segment $[x,y_{x}]$ is parallel to the hyperplane $\R^{d-1}\times\set{0}$ and has length 
	\begin{equation}\label{eq:dist_x-y_x}
	\enorm{y_{x}-x}\leq \diam\br{Q}=s^{2}.
	\end{equation}
	Therefore whenever $x,x'\in X$ and $\proj_{\R^{d-1}}\mf{U}_{x}\cap \proj_{\R^{d-1}}{\mf{U}_{x'}}\neq \emptyset$ we have 
	\begin{equation*}
	\enorm{\proj_{\R^{d-1}}(x')-\proj_{\R^{d-1}}(x)}\leq 2s^{2}+ r_{x}+r_{x'}\leq 3s^{2},
	\end{equation*}
	where the last inequality comes from imposing the condition $r_{x}(s)\leq s^{2}/2$ for every $x\in X$. Therefore, whenever $x$ and $x'$ are distinct points of the $s$-separated set $X$ and $\proj_{\R^{d-1}}\mf{U}_{x}\cap \proj_{\R^{d-1}}{\mf{U}_{x'}}\neq \emptyset$ we have
	\begin{equation*}
	\abs{\proj_{d}(x')-\proj_{d}(x)}\geq s-3s^{2}\geq s^{2}\geq 2\sup_{w\in X}r_{w}(s).
	\end{equation*}
	We deduce that the collection $\br{\mf{U}_{x}}_{x\in X}$ is pairwise disjoint. Moreover, the condition $X\subseteq \R^{d-1}\times \sqbr{\frac{1}{2}+s,H+\frac{1}{2}-s}$ and the fact that each $[x,y_{x}]$ is parallel to $\R^{d-1}\times \set{0}$ ensure that each set $\mf{U}_{x}$ in this collection is contained in $\R^{d-1}\times \br{\frac{1}{2},H+\frac{1}{2}}$. 
	
	
	Let $\Xi\colon \R^{d}\to\R^{d}$ be the $16$-bilipschitz mapping given by the application of Theorem~\ref{thm:simultaneous_switching} to $X$, $(y_{x})_{x\in X}$ and $(r_{x})_{x\in X}$.
	Here we need one last condition on $r_{x}(s)$, namely $r_{x}(s)\leq \enorm{y_{x}-x}/2$ for all $x\in X$ with $y_{x}\neq x$; the choice
	\begin{equation*}
	r_{x}(s):=\begin{cases}
	\frac{1}{2}\enorm{y_{x}-x} & \text{if $y_{x}\neq x$,}\\
	\frac{1}{2}s^{2} & \text{otherwise}
	\end{cases}
	\end{equation*}
	satisfies both conditions we required. The properties \eqref{Psi_i}--\eqref{Psi_iv} of $\Xi$ are now easily verified.
	
	\paragraph{Construction of $\Omega$.} 
	We construct a $16N^{2}H^{2}$-bilipschitz mapping $\Omega$ with the following properties:
	\begin{enumerate}[(I)]
		\item\label{rho_i} $\Omega(x)=x$ for all $x\in\R^{d-1}\times \br{\R\setminus \br{\frac{1}{2},H+\frac{1}{2}}}$.
		\item\label{rho_ii} $\proj_{\R^{d-1}}\circ \Omega\circ \Xi(x)=\proj_{\R^{d-1}}\circ \Xi(x)$ for all $x\in X$.
		\item\label{rho_iii} $\proj_{d}\circ \Omega\circ \Xi(x)=m$ for all $x\in X_{m}$ and $m=1,2,\ldots,H$.
	\end{enumerate}
	For each $m\in [H]$ and $x\in X_{m}$ let 
	\begin{align*}
	v_{x}&:=\Xi(x), & y_{x}:=(\proj_{\R^{d-1}}\circ \Xi(x),m)\in \R^{d},\\
	r_{x}&:=\begin{cases}
	\min\set{\frac{1}{2N},\frac{\enorm{y_{x}-v_{x}}}{2}} & \text{if }v_x\neq y_x,\\
	\frac{1}{2N} & \text{if }v_x=y_x,
	\end{cases} & \mf{U}_{x}:=B([v_{x},y_{x}], r_{x}).
	\end{align*}
	The choice of these objects, together with property \eqref{Psi_iii} of $\Xi$, ensures that the collection $(\mf{U}_{x})_{x\in X}$ is pairwise disjoint. Moreover, the hypotheses on $X$, the bounds $r_{x}<1/N<s$ and \eqref{Psi_iv} ensure that each set $\mf{U}_{x}$ with $x\in X$ is contained in $\R^{d-1}\times \br{\frac{1}{2},H+\frac{1}{2}}$. 
	
	
	Let the $16N^{2}H^{2}$-bilipschitz mapping $\Omega\colon \R^{d}\to \R^{d}$ be given by Theorem~\ref{thm:simultaneous_switching} applied to $\set{v_{x}\colon x\in X}$ in place of $X$, $v_{x}$ in the role of $x$, $(y_{x})_{x\in X}$ and $(r_{x})_{x\in X}$. The properties \eqref{rho_i}--\eqref{rho_iii} of $\Omega$ are clear from the construction. 
	\paragraph{Properties of $\Psi:=\Omega\circ \Xi$.}
	Setting $\Psi:=\Omega\circ \Xi$ we get that $\Psi$ is $2^{8}N^{2}H^{2}$-bilipschitz. The properties \eqref{Xi_i}--\eqref{Xi_iii} of $\Psi$ follow easily from \eqref{Psi_i}--\eqref{Psi_iv} for $\Xi$ and \eqref{rho_i}--\eqref{rho_iii} for $\Omega$.
\end{proof}

\end{document}
